\section{Discussion}
\label{sec:discussion}

\subsection{Why Subjective Logic Coordination Earns Its Place}
A superficial reading of raw predictive accuracy might view the $+0.0084$ F1-Macro gain of AEGIS-SL ($0.7190$) over monolithic fusion ($0.7106$) as modest. However, our 2$\times$2 mechanistic ablation (§\ref{sec:results}) reveals that the Subjective Logic opinion tensor contributes $+0.0072$ F1 while dynamic Brier credibility tracking adds $+0.0048$ F1. 

More importantly, raw F1 does not reflect the load-bearing architectural advantages of AEGIS-SL:
\begin{enumerate}
    \item \textbf{Calibrated Uncertainty}: AEGIS-SL achieves an ECE of $0.0925$ and mean epistemic uncertainty $u=0.0957$. Monolithic models (BL2) provide zero uncertainty bounds.
    \item \textbf{Missing-Modality Safety}: Under sensor failure, AEGIS-SL exhibits strictly monotonic uncertainty growth (Table~\ref{tab:monotonicity}), whereas monolithic models crash or require uncalibrated zero-filling.
    \item \textbf{Auditability \& Provenance}: AEGIS-SL generates an immutable log detailing per-agent belief contributions, enabling interactive GUI rendering.
\end{enumerate}

\subsection{The BL3 Failure Pattern as a General Result}
The decoupling of AUROC ($0.7127$) and F1-Macro ($0.2256$) in Baseline 3 represents a fundamental failure mode of LLM-as-arbitrator architectures. When agents emit conflicting sensory inputs, LLM prompt aggregation collapses onto the majority class ('Dry'), creating severe class-imbalance failure. This result demonstrates that text-based prompt arbitration is unsuited for safety-critical environmental risk monitoring.

\subsection{Limitations}
This study has three primary limitations:
\begin{enumerate}
    \item \textbf{Single AOI Event}: Evaluation is focused on the Brisbane 2022 event ($200$ cells $\times 24$ days). Spatial and temporal generalization across global flood regimes requires further empirical testing.
    \item \textbf{Proxy Cohort Size}: The $N=12$ analyst study is directionally supportive ($\Delta\text{Likert} = +0.82$) but statistically underpowered for large-scale usability conclusions.
    \item \textbf{Hybrid Head Dependency}: The final prediction head relies on LightGBM decision trees rather than propagating Dirichlet evidence end-to-end.
\end{enumerate}

\subsection{Broader Applicability}
The AEGIS 5-agent evidential architecture is domain-agnostic and directly transferable to other multi-modal climate disaster domains:
\begin{itemize}
    \item \textbf{Wildfire Risk Monitoring}: Fusing thermal infrared, fuel moisture, wind vectors, terrain slope, and satellite active fires.
    \item \textbf{Deforestation Event Detection}: Fusing optical canopy loss, SAR radar backscatter, road proximity, and indigenous land boundary text notices.
    \item \textbf{Urban Heat Island \& Air Quality Events}: Fusing thermal sensors, PM2.5 monitoring, traffic densities, and atmospheric inversion warnings.
\end{itemize}

\subsection{Future Work}
Future research will explore three directions:
\begin{enumerate}
    \item \textbf{Multi-Event Generalization}: Evaluating AEGIS across global benchmark floods (e.g., Auckland January 2023 and Dubai April 2024).
    \item \textbf{Generative Dirichlet SL Head}: Replacing the LightGBM hybrid head with a fully evidential neural network architecture propagating Dirichlet distributions end-to-end.
    \item \textbf{Human-in-the-Loop Credibility Seeding}: Allowing domain hydrologists to seed agent credibility priors $\gamma_i$ interactively during unfolding emergencies.
\end{enumerate}
